% ──────────────────────────────────────────────────────────────────────────────
%  Section 6 — Parametric sweeps of the peak-to-peak spread
% ──────────────────────────────────────────────────────────────────────────────

\section{Parametric sweeps of $\Delta_{\max}$}
\label{sec:sweeps}

\autoref{eq:Deltamax} is evaluated on two single-variable sweeps of the
$\SI{110}{\kilo\volt}$ asymmetric configuration of
\autoref{sec:example}. \autoref{fig:sweep-S} fixes the line length at
$\ell=\SI{25}{\kilo\metre}$ and sweeps $|S|$ from no load up to
$\SI{200}{\mega\volt\ampere}$ along the direction of the example
operating point, that is $\arg(S)=\arg(80+\jj 10\,\si{\mega\volt\ampere})
\approx 7.13\si{\degree}$. \autoref{fig:sweep-length} fixes the apparent
power at the example value $S=80+\jj 10\,\si{\mega\volt\ampere}$ and
sweeps the line length from $\SI{1}{\kilo\metre}$ up to
$\SI{100}{\kilo\metre}$. The filled marker on each curve is the common
anchor at $S=80+\jj 10\,\si{\mega\volt\ampere}$ and
$\ell=\SI{25}{\kilo\metre}$.

\begin{figure}[ht]
\centering
\begin{tikzpicture}
  \begin{axis}[
    width=0.95\linewidth,
    height=7cm,
    xlabel={$|S|$ \,[MVA] at $\arg S = \arg(80{+}\jj 10)$},
    ylabel={$\Delta_{\max}$ \,[A]},
    xmin=0, xmax=200,
    ymin=0,
    grid=both,
    minor tick num=1,
    tick label style={/pgf/number format/fixed},
  ]
    \addplot+[
      no marks, line width=0.9pt, color=primary,
    ] table[x=S_MVA, y=spread_A] {data/sweep_S.dat};

    \addplot[only marks, mark=*, mark size=2.4pt, color=primary,
             forget plot]
      coordinates {(80.6226, 38.4665)};
  \end{axis}
\end{tikzpicture}
\caption{Peak-to-peak modular spread $\Delta_{\max}$ at the local
terminal versus apparent-power magnitude $|S|$, with the line length
held at $\ell=\SI{25}{\kilo\metre}$ and the loading direction held at
the operating power factor of the example. The filled marker is the
anchor at $S=80+\jj 10\,\si{\mega\volt\ampere}$.}
\label{fig:sweep-S}
\end{figure}

\autoref{fig:sweep-S} confirms the heavy-load asymptote of
\autoref{eq:heavyloading}: the spread is linear in $|S|$ across the
plotted range, with a slope set by the spread of $\sqrt{\xi_{k}(\ell)}$
across the three phases. The line passes through the origin to within
the no-load offset $\max_{k}|b_{k}|-\min_{k}|b_{k}|<\SI{0.05}{\ampere}$
that becomes visible only at $|S|\lesssim\SI{0.5}{\mega\volt\ampere}$
and is not separable from zero at the scale of the figure.

\begin{figure}[ht]
\centering
\begin{tikzpicture}
  \begin{axis}[
    width=0.95\linewidth,
    height=7cm,
    xlabel={Line length $\ell$ \,[km]},
    ylabel={$\Delta_{\max}$ \,[A]},
    xmin=0, xmax=100,
    ymin=38, ymax=39,
    grid=both,
    minor tick num=1,
    tick label style={/pgf/number format/fixed},
  ]
    \addplot+[
      no marks, line width=0.9pt, color=primary,
    ] table[x=length_km, y=spread_A] {data/sweep_length.dat};

    \addplot[only marks, mark=*, mark size=2.4pt, color=primary,
             forget plot]
      coordinates {(25, 38.4665)};
  \end{axis}
\end{tikzpicture}
\caption{Peak-to-peak modular spread $\Delta_{\max}$ at the local
terminal versus line length $\ell$, with the apparent power held at the
example value $S=80+\jj 10\,\si{\mega\volt\ampere}$. The filled marker
is the anchor at $\ell=\SI{25}{\kilo\metre}$.}
\label{fig:sweep-length}
\end{figure}

\autoref{fig:sweep-length} shows that for fixed $S$ the spread is
remarkably flat across the entire span of practical
$\SI{110}{\kilo\volt}$ corridor lengths. The variation across
$\ell\in[1,\,100]\,\si{\kilo\metre}$ is bounded by
$\SI{0.2}{\ampere}$ around a base value of roughly
$\SI{38.4}{\ampere}$, which is at the resolution edge of substation
ammeters and consistent with the observation that, for fixed loading,
the unbalance pattern is set by the asymmetry of the per-unit-length
matrices and not by the length. The length enters $\Delta_{\max}$
through the ABCD blocks of $\matM(\ell)=\ee^{\matF\ell}$, but for the
geometry under study the resulting variation of
$\sqrt{\xi_{k}(\ell)}$ across the three phases is essentially
invariant in $\ell$; the asymptote of \autoref{eq:heavyloading} is
attained almost immediately and stays there as the line is lengthened.
